\todo{Chương 2 chỉ có khung: nhóm tự viết văn bản từ các ý chính bên dưới và trích nguồn đã đọc; các mục tương ứng nội dung môn học (chương 6, 8).}

\section{Tổng quan về kinh doanh thông minh}
\begin{itemize}
  \item Khái niệm BI, vòng đời từ dữ liệu đến quyết định; BI và DW bổ trợ nhau \cite{kimball2013}.
  \item Các cấp phân tích: mô tả, chẩn đoán, dự đoán, đề xuất; phạm vi bài tập (mô tả và chẩn đoán).
\end{itemize}

\section{Kiến trúc và công cụ BI}
\begin{itemize}
  \item Các thành phần: nguồn dữ liệu, kho / data mart, mô hình ngữ nghĩa, báo cáo và dashboard, phân phối.
  \item So sánh các công cụ phổ biến (Power BI, Tableau, Qlik) và lý do chọn Power BI \cite{powerbi_docs}.
  \item Chế độ Import so với DirectQuery; Publish to web và các lưu ý về chia sẻ dữ liệu.
\end{itemize}

\section{KPI và đo lường hiệu quả}
\begin{itemize}
  \item KPI, chỉ số dẫn dắt và chỉ số kết quả; tiêu chí SMART; cây KPI.
  \item Liên hệ KPI của bài toán: GMV, AOV, tỷ lệ giao đúng hạn, điểm review, tỷ lệ chuyển đổi MQL sang seller (mục 3.3.5).
\end{itemize}

\section{Trực quan hóa dữ liệu và thiết kế dashboard}
\begin{itemize}
  \item Nguyên tắc chọn biểu đồ theo mục đích (so sánh, xu hướng, cấu thành, phân bố, quan hệ, bản đồ).
  \item Nguyên tắc bố cục, màu sắc, định dạng có điều kiện, slicer và drill-down; các loại dashboard (theo thời gian, OLAP, vận hành, thống kê).
\end{itemize}

\section{Mô hình ngữ nghĩa và ngôn ngữ DAX}
\begin{itemize}
  \item Bảng, quan hệ, cardinality, role-playing date và \code{USERELATIONSHIP}; ngữ cảnh bộ lọc và ngữ cảnh hàng \cite{dax_docs}.
  \item Measure so với cột tính toán; time intelligence (MoM, YoY).
\end{itemize}

\section{Xu hướng}
\begin{itemize}
  \item BI tự phục vụ, BI nhúng, tích hợp AI; dữ liệu lớn (giáo trình chương 8).
\end{itemize}
